\phantomsection\label{first-target-finite-dimensional-trace}

Autoformalized; awaiting human review. All three section companions and their shared mathematics have compiled with Lean 4.21.0-rc3 and the pinned Mathlib revision. The chapter-wide build and axiom audit passed most recently on 2026-11-05 for 297 logical declarations, including generated helpers. The evidence record is \texttt{lean/audits/chapter00-2026-11-05-editor/record.json}; the recorded source hashes identify the version that passed those checks. Steven's review of the correspondence between the manuscript and the formal statements remains pending independently.

Section 0.1 formalizes rectangular cyclicity, similarity invariance, basis independence of endomorphism trace, and the trace as the sum of characteristic roots with algebraic multiplicity. Empty matrices are included, and the last theorem does not assume diagonalizability. The concrete examples include the Jordan matrix, its eigenspace dimension, and the absence of an eigenbasis. Mathlib supplies the standard matrix, determinant, polynomial, basis, and trace constructions, including the proved splitting results for complex polynomials.

\phantomsection\label{chapter-zero-formalization-roadmap}

Section 0.2 formalizes finite orthonormal coefficient extraction and the orthogonal remainder, constructs the finite-rank operator and its kernel, and proves the integral representation and diagonal trace identity. The operator trace is compared across finite-dimensional supporting subspaces. The spectral specialization, counting-measure interpretation, and normalized interval rank-one example are included. Representative changes are handled explicitly: changing the finitely many basis representatives preserves the diagonal integral, while the zero and diagonal-indicator kernels demonstrate the limitation of arbitrary product-null kernel changes.

Section 0.3 formalizes the relative kernel's integrability, its period expansion and operator pairing, and the diagonal specialization. It also constructs bounded multiplication by an $L^\infty$ function on $L^2$, proves the weighted integral and compressed-trace identity, and derives its diagonal specialization. The two-point example uses counting measure and includes its ordinary, relative, and weighted measurements.

The shared library reuses Mathlib's linear algebra and measure theory. The book's inner product is linear in its first argument; the formalization reverses Mathlib's inner-product arguments to implement that convention. Since $L^2$ elements are equivalence classes, the integral formula for an operator's action is stated almost everywhere. The finite-sum kernel representatives used for diagonal integration are specified separately. The section comparisons explain these conventions and any different proof route for Steven to review.

Run \texttt{lake build} from the \texttt{traceformulas-vscode/} project directory to compile the shared library and all three companions. The Lake configuration is in that common parent of the shared library and manuscript sources, so the editor can find the imports when opening a section's Lean file. From the same project directory, \texttt{python3 scripts/audit\_lean.py} builds the project and records an axiom audit against the saved source snapshot. The gate examines kernel dependency closures and permits only Lean's ordinary foundational axioms; it does not approve the mathematical correspondence on Steven's behalf.

The Compare tab reads each section's \texttt{text.tex} and \texttt{Formalization.lean} directly, together with \texttt{compare.tex}. The website displays those sources but does not execute Lean or certify the proofs. Steven explicitly requested the existing chapter-level video script before correspondence review. It remains an unreviewed draft; no recording or production approval is implied.
